% !TEX root = main.tex
\section{Related Work (Cont.)}\label{sec:related-work-cont}

\paragraph{Autoformalization.} Complementing this line of work, \citet{wang2024theoremLlama} present \textsc{TheoremLlama}, a framework aimed at enhancing LLM translation into Lean 4.
Drawing on over 100K proof examples from the Mathlib4 library, TheoremLlama employs a novel natural-language-to-formal-language (NL-FL) bootstrapping strategy and iterative proof synthesis.
This enables the reuse of verified examples as templates for future translations.
The framework achieves 36.48\% and 33.61\% accuracy on the MiniF2F-Valid and MiniF2F-Test benchmarks, respectively—surpassing GPT-4 by more than ten percentage points on both.

\textbf{Techniques for Code Verification.} \textsc{ImProver}~\citep{ahuja2024improver} introduces a Lean-aware Chain-of-States prompting loop that integrates retrieval, best-of-$n$ sampling, and iterative correction to rewrite formal proofs with improved properties.
By optimizing for metrics such as brevity and readability, ImProver reduces the number of tactics by half, doubles proof readability, and boosts theorem prover acceptance rates by over 80\%.
In addition, \textsc{Clover}~\citep{sun2024clover} implements a closed-loop pipeline in which an LLM first generates code, docstrings, and formal annotations, then uses reconstruction-based prompting to enforce consistency across these outputs, and finally applies SMT-based verification to validate correctness.
Evaluated on the CloverBench suite of Dafny programs, Clover accepts 87\% of correct solutions, rejects 100\% of flawed ones, and even uncovers bugs in human-written code—demonstrating the power of hybrid generation-verification pipelines.
In a complementary direction, \citet{zhou2025graders} improve general-purpose LLM-based graders by augmenting them with “privileged” information such as gold-standard solutions, grading rubrics, and detailed annotations.
When necessary, the system provides targeted hints back to candidate models.
This approach achieves grading performance on par with or exceeding that of specialized systems—and even expert humans—on difficult programming benchmarks.

\textbf{Agentic Frameworks and Tools for Code Verification.} \textsc{Trace}~\citep{cheng2024trace} proposes \emph{generative optimization}, tuning entire computational workflows—including code, prompts, tool calls, and error signals—by treating execution traces as gradients in the OPTO framework.
With a PyTorch-like API and the LLM-based optimizer OptoPrime, it supports diverse tasks such as prompt tuning, debugging, and robot control, rivaling specialized optimizers.
Building on modular composition, \textsc{DSPy}~\citep{khattab2024dspy} treats LLM calls as declarative modules in a computational graph.
Users define concise input–output signatures, and DSPy’s compiler automatically bootstraps or fine-tunes pipelines using built-in “teleprompters.”
This enables a few-line programs to outperform expert-crafted prompts in math, QA, and agent workflows.
Extending agentic capabilities to formal reasoning, \textsc{Pantograph}~\citep{aniva2025pantograph} offers a programmatic interface to Lean 4 with support for advanced proof search.
It exposes internal proof states and tactics for integration with learning agents, replacing human-facing interfaces with API-level control.

\subsection{Related Work: Scalable Oversight}

\paragraph{Scalable Oversight with LLM Judges}
Researchers have proposed multi-agent oversight frameworks such as debate and consultancy, where weaker LLMs act as judges supervising stronger agents \citep{kenton2024oversight,Chan2024ChatEval}.
Debate, in which two models argue opposing sides before a judge, consistently improves truth-tracking compared to single-pass QA.
\citet{kenton2024oversight} demonstrate that debate reduces the likelihood of judges being misled in math and code tasks, highlighting it as a scalable oversight mechanism.

\paragraph{Self-Improving Critics.}
In the absence of stronger supervisors, models can bootstrap judgment skill through self-critique.
The Self-Evolving Critic (SCRIT) framework trains LLMs to generate critiques of flawed outputs, validate corrections, and iteratively refine their ability to spot errors.
\citet{zheng2024scrit} report improvements of over 10\% on critique benchmarks, showing that self-critique fosters scalable, autonomous oversight.

\paragraph{Certifying Judges.}
Naïve LLM judges are often inconsistent, rewarding stylistic cues over correctness \citep{Kocmi2024HumanVsLLM,Li2024TrustLLMJudgments}.

\paragraph{Validity-Centered Evaluation.}
Certifying judges requires not only behavioral properties but also principled validity evidence.
\citet{salaudeen2025measurement} introduce a framework that distinguishes between measurement, claims, and evidence, grounded in psychometrics.
They emphasize construct, content, and criterion validity, and systematically enumerate threats to validity.
This framework provides a blueprint for reporting and interpreting judge certification results, complementing identity/monotonicity checks with broader validity cards.

\paragraph{Robust Evaluation without Ground Truth.}
Standard judge queries can be brittle in the absence of ground truth, particularly under adversarial manipulation \citep{Zhuo2024ICEScore,TongZhang2024CodeJudge}.
\citet{robertson2025stepbystep} model the overseer as an agent and propose using information-theoretic, incentive-compatible queries instead of direct ratings.
They introduce bounded $f$-divergence measures such as TVD-MI, which exhibit polynomial robustness to adversarial strategies, unlike unbounded KL-based criteria.
This work highlights a pathway toward robust, certifiable judges when explicit ground truth is unavailable.

\paragraph{Single-judge certification vs.\ peer-prediction consensus.}
Information-elicitation approaches such as TVD-MI \citep{robertson2025stepbystep} certify trust by asking multiple judges whether their outputs share information about the same source---fundamentally a consensus-based protocol over a judge population.
Property certification is complementary: identity and monotonicity properties evaluate a single judge in isolation, with no agent-population dependency.
This makes property certification deployable in settings where consensus is unavailable (single-judge inference, novel domains, sparse model panels), and our human-agreement assessment protocol amortizes the human-annotation cost that consensus mechanisms also incur.

\paragraph{Consensus aggregation vs.\ property certification.}
\citet{denisovblanch2026consensus} stress-test five inference-time aggregation strategies---majority vote, highest-confidence selection, confidence- and prediction-weighted votes, and the Surprisingly Popular algorithm---across multiple benchmarks (HLE, BoolQ, Com2Sense, and Predict-the-Future) and open-source models spanning roughly 4B to 235B parameters; even at $25\times$ the cost of single-sample inference no aggregator consistently beats a single sample in verifier-absent domains, and several amplify shared misconceptions.
They trace the failure to structurally correlated LLM errors: in MATH plurality-wrong cases, the most common incorrect answer accounts for 65--87\% of wrong responses, and in an out-of-distribution random-ASCII control with no ground truth, different models still agree at Cohen's $\kappa$ as high as ${\approx}0.35$, indicating shared inductive biases rather than shared knowledge.
A central diagnosis is that self-reported confidence and consensus track \emph{what the crowd will say} more reliably than \emph{what is true}, separating social prediction from truth verification, and the paper concludes that inference-time scaling improves truthfulness only when an external verifier is available.
Our work sits on the verifier-present side of that boundary: instead of aggregating across a judge population, we evaluate a single judge against falsifiable properties (P1--P4) over Lean~4-checked gold artifacts and explicitly constructed bug/specification ladders, yielding a metric for assessing appropriate reliance with no inter-judge consensus dependency.
For oracle-free extensions, TVD-MI-style cross-judge information should be used as a robustness diagnostic rather than as a truth proxy, side-stepping the consensus-as-truth failure mode they document; conversely, their negative results sharpen the case for property-based, single-judge certification whenever a verifier or reference checked against known outcomes exists.

\paragraph{Specializing judges to a test set vs.\ certifying judges across configurations.}
\citet{finkelstein2024specialist} propose the \emph{Specialist} method, which deliberately abandons cross-test-set generalization and instead overfits a prompted LLM autorater to a single fixed benchmark.
For each test example, in-context demonstrations are drawn from historical human ratings of the \emph{same source}, yielding 54--119\% character-level F1 gains over \textsc{XCOMET} on the WMT'23/'24 MQM tasks.
Their work and ours share a common motivation---amortizing the expensive part of human evaluation rather than re-annotating from scratch---but amortize along orthogonal axes.
Specialist amortizes \emph{within} a test set: it reuses many historical ratings (under a pseudo-SxS, single-rater-per-input regime) as inference-time demonstrations, so the judge must be re-specialized for each new benchmark and, as the authors show, to the individual rater.
Property certification instead amortizes \emph{across judge configurations}: a one-time human-agreement assessment panel relates falsifiable behavioral properties (P1--P4) to human agreement, after which any new $(\text{model},\text{prompt},n_\text{repeats},\text{aggregation})$ tuple is screened by fully automated perturbation checks, with no human-vs-gold labels required on the perturbed variants.
The two approaches are complementary: Specialist still consumes human ratings at scoring time and presumes a curated pool of same-source annotations, whereas our property-certification protocol targets the setting where new judges proliferate faster than any per-test-set rating pool can be maintained.
Notably, their finding that a Specialist autorater specializes to a \emph{single} rater echoes our own concern with human-pass dispersion; both works suggest single-rater ground truth is an imperfect ceiling, motivating the multi-rater held-out protocol we defer to future work.

\paragraph{Formal Methods as Oversight.}
Proof assistants provide binary, kernel-level guarantees.
\citet{zhou2024donttrustverify} autoformalize math solutions into Isabelle/HOL and show that filtering via proof checking outperforms majority voting by 12\% on GSM8K.
\citet{liu2025safe} extend this approach with SAFE, which verifies every intermediate step in Lean 4 and introduces the FormalStep benchmark, reducing hallucinations and enabling fine-grained oversight.

\paragraph{Agentic Iterative Pipelines.}
\citet{ospanov2023apollo} integrate Lean feedback into recursive repair cycles, boosting general-purpose LLM proof success from 3--7\% to over 40\% on miniF2F, and achieving 75\% with fewer than 1000 samples.
\citet{miranda2026veribench} extend this paradigm to program verification, showing that self-debugging agents quadruple success compared to one-shot prompting.
Incorporating an LLM judge improved completeness on security tasks but required tuning to avoid regressions on simpler cases.

\paragraph{Benchmarks for Judge Reliability.}
Benchmarks such as VeriBench and FormalStep quantify judge reliability across compilation, proof success, and specification quality.
\citet{miranda2026veribench} report normalized specification scores of 0.470 for one-shot prompting versus 0.615 for self-debugging agents.
Outside formal mathematics, MT-Bench and G-Eval provide complementary judge-evaluation settings using pairwise preferences and rubric-guided scoring, underscoring the need for human-agreement analysis and validity reporting \citep{Zheng2023MTBench,Liu2023GEval}.
Together with validity frameworks and robust information-theoretic scoring, these benchmarks advance the agenda of certifying LLM judges for formal verification.

% https://gemini.google.com/app/d39eba54ee6763e1
\paragraph{Evaluating Autoformalization with Reference-Free Judges.}
In a complementary line of work, \citet{zhang2025goldstandardsepistemicensemble} propose an ensemble of LLM judges for the autoformalization task, using a fine-grained taxonomy to score formalizations without a gold-standard reference.
The primary strength of this approach is its reference-free nature, offering a scalable evaluation pathway when gold standards are unavailable or prohibitively expensive to create.
However, the work's motivation for abandoning gold-standard comparison is based on limited evidence: a small, non-random sample of "potentially incorrect" formalizations pre-selected by an LLM, whose own reliability is uncertified.
Furthermore, the framework validates its judge by correlating with human preference scores—a subjective proxy—rather than the more objective, machine-checkable goal of logical equivalence to a known-correct specification.
Critically, even if minor errors exist in gold-standard datasets, recent work has shown that such noise rarely alters the relative ranking of models, as errors are typically correlated across systems \citep{salaudeen2024imagenot}.
This suggests that a small number of flawed examples does not invalidate the utility of an entire benchmark for comparative evaluation.
Our work, in contrast, proceeds from gold references that Lean kernel-check relative to explicitly stated predicates, and emphasizes \emph{validation-set evidence} for pairwise judge scores---the preferred rubric-aligned target when verifier-backed references exist.

% https://chatgpt.com/c/68dc36c7-c4b0-8327-9dcd-694c10fa5d93?model=gpt-5-pro
\subsection{HCI for Solving Scalable Oversight}

\paragraph{From more evaluation to \emph{human-centered} oversight.}
HCI reframes scalable oversight as a problem of matching reliance to demonstrated reliability: design interactions that let people verify, steer, and place \emph{appropriate} trust in LLM outputs rather than maximizing raw trust \citep{bach2022trusthci}. Interface choices (e.g., citations, explanations, confidence cues) can substantially shift user trust---sometimes in misleading ways---so the goal is to expose process and embed lightweight verification instead of presenting one-shot answers \citep{ding2025citations}. User studies in AI pair-programming similarly surface needs to communicate model performance, provide control over output style/constraints, and expose uncertainty to avoid over-reliance \citep{wang2024trustcode}.%
% :contentReference[oaicite:0]{index=0}

\paragraph{Hybrid human$\leftrightarrow$AI evaluation.}
HCI work operationalizes scalable evaluation by combining expert/lay judgments with automated judges. Humans annotate failures or verify slices of output; systems then train or condition LLM-based ``reference evaluators'' on those signals to scale up assessment while staying anchored to what users value (accuracy, safety, usefulness). This human-in-the-loop pattern complements LLM-as-judge approaches---whose reliability varies by task and setup \citep{Kocmi2024HumanVsLLM,Li2024TrustLLMJudgments}---and is especially relevant when pass@k or single-shot metrics are brittle.%
% :contentReference[oaicite:1]{index=1}

\paragraph{Interfaces that \emph{operationalize} oversight.}
HCI systems instantiate oversight as interaction loops that make correctness checkable \emph{during} generation: (i) \emph{test-driven code generation} (TiCoder) elicits yes/no judgments about candidate unit tests before showing code; a few interactions reduced cognitive effort and lifted pass@1 by $\sim$45.7\% on coding benchmarks \citep{naik2024ticoder}; (ii) \emph{explain$\rightarrow$correct loops} (IntelliExplain) pair each snippet with a natural-language rationale that users can critique, improving novice task success by $\sim$11--25\% \citep{yan2024intelliexplain}; (iii) \emph{pair debugging} (HypoCompass) has humans hypothesize faults while the LLM attempts fixes, yielding $\sim$12\% learning gains \citep{ma2024hypocompass}; (iv) \emph{multi-step pipelines} (PromptChainer) expose intermediate steps (e.g., tests $\rightarrow$ verification $\rightarrow$ proof attempt), enabling user intervention where failures arise \citep{wu2022promptchainer}; and (v) \emph{source-grounded verification} (SymGen) highlights claim-level spans backed by exact sources and leaves unsupported text unhighlighted, speeding verification by $\sim$20\% and raising user confidence \citep{shen2024symgen}. Collectively, these patterns shift trust from fluency to \emph{verifiable progress}.%
% :contentReference[oaicite:2]{index=2}

\paragraph{Applying HCI patterns to formal specs/theorems.}
Formal kernels (Lean/Isabelle) already deliver binary checks; HCI contributes the \emph{workflow glue} that helps users evaluate \emph{which} theorem/spec to prove and whether it is complete. Two actionable patterns: (1) \emph{counterfactual spec probes}---UI toggles that remove a postcondition or inject a known bug and display the impact on tests/judge score---help detect missing clauses or overfitting; and (2) \emph{stepwise proof oversight}---surfacing intermediate goals and marking unverified suggestions---keeps the human in control of proof search. In our setting, these patterns interface naturally with a graded, semantics-sensitive evaluator.%
% :contentReference[oaicite:3]{index=3}

% \paragraph{Complementarity with a \emph{certified} formal judge.}
% % :contentReference[oaicite:4]{index=4}

% \paragraph{Design takeaways for scalable oversight.}
% \emph{Expose process}, not just outputs (\citealp{wu2022promptchainer}); \emph{facilitate incremental verification} (\citealp{naik2024ticoder,shen2024symgen}); \emph{empower corrective feedback} (\citealp{yan2024intelliexplain}); and \emph{communicate uncertainty honestly} (\citealp{wang2024trustcode,ding2025citations}). These HCI patterns consistently reduce cognitive load, improve objective accuracy, and---critically---align user trust with actual reliability across code/spec/theorem tasks.%
% % :contentReference[oaicite:5]{index=5}

